\documentclass[12pt]{article}
\usepackage[utf8]{inputenc}
\usepackage[T1]{fontenc}
\usepackage{hyperref}
\title{Purisime\~no Language Resources for K-12 Curriculum Development}
\author{Compiled by Claude Sonnet 5.5 (claude-sonnet-5-5) for Andr\'as Kornai}
\date{}
\begin{document}
\maketitle

\section{Introduction}
Purisime\~no (ISO 639-3 \texttt{puy}) is a dormant Chumashan language of California, once spoken around Mission La Pur\'isima Concepci\'on. It has no first-language speakers. The documentation is largely John P.\ Harrington's field notes (about 2700 slides according to search results) and older word lists. No school programme was found.

\section{Primary Educational Resources}
No primers, textbooks or school curricula were found. The lexicon and grammar below are technical works, though the lexicon is described as usable by non-specialists.

\section{Linguistic and Reference Materials}
\begin{itemize}
\item Grammar and Lexicon of Purisime\~no Chumash (T.\ Henry-Rodriguez), University of Hawai\textquoteleft i ScholarSpace collection: a Purisime\~no--English and English--Purisime\~no lexicon (2021) and a grammar, based on Harrington's data. \url{https://scholarspace.manoa.hawaii.edu/collections/e14cc0dd-2be8-4264-962d-01652abd75e9}
\item A ScholarSpace file from the same search results; I did not open it, so its exact content is unverified. \url{https://scholarspace.manoa.hawaii.edu/server/api/core/bitstreams/e821a88f-71a7-4040-8f73-9c9802f2babb/content}
\item Wikipedia, Purisime\~no language. \url{https://en.wikipedia.org/wiki/Purisime%C3%B1o_language}
\item UC Berkeley Survey of California and Other Indian Languages, Purisime\~no page. \url{https://cla.berkeley.edu/languages/purisimeno.html}
\item UBC wiki: California Indian Linguistic Records, the Mission Indian vocabularies of Alphonse Pinart. \url{https://wiki.ubc.ca/Documentation:RelLex/California_Indian_Linguistic_Records:_The_Mission_Indian_Vocabularies_of_Alphonse_Pinart}
\end{itemize}

\section{Religious and Community Resources}
No scripture or community organisation specific to Purisime\~no was found in search results.

\section{Digital and Audio Resources}
\begin{itemize}
\item The family overview on internationalphoneticalphabet.org, ``the family of Chumash languages''. \url{https://www.internationalphoneticalphabet.org/ineseno-chumash/the-family-of-chumash-languages/}
\end{itemize}
No audio was found. Search results mention the Purisime\~no Project of the Western Institute for Endangered Language Documentation (wieldoc.org), but I did not see a direct URL in the results, so none is given.

\section{Curriculum Development Context}
A heritage curriculum would rest on the Henry-Rodriguez lexicon and grammar, together with the Harrington material. Because the language is dormant, any use in K-12 should be developed with the descendant Chumash community. Related Chumash revitalization work may offer models.

\end{document}